% ── Author's Note ─────────────────────────────────────────────

\chapter*{A Note from the Author}
\addcontentsline{toc}{chapter}{A Note from the Author}

\thispagestyle{plain}

This novel is a first attempt at something that may require several more.

The four figures at its centre---Guyon, Mackandal, Payne-Gaposchkin,
and Julius Caesar Scaliger---are not equally historical. The first
three are reconstructions, drawn from the record but interpreted
through a framework the record did not provide. The fourth is
speculation entirely: a figure placed in a future that does not
exist in order to suggest that what the first three discovered
independently is not finished being discovered.

The framework behind the novel concerns the different ways that
systems---minds, networks, instruments, practices---move from
multiplicity toward coherence, or fail to. Guyon called this
process prayer. Mackandal encoded it in objects that circulated
through a network. Payne extracted it from light. Julius finds
it by moving through water until the water tells him where the
problem is. The novel argues, without stating, that these are
versions of the same thing.

A word about the dialogue, or its near-absence.

The characters in this novel speak less than characters in most
novels speak, and when they speak, the speech is often indirect,
reported, or absorbed into the surrounding prose rather than
rendered as direct exchange. This was partly intentional and
partly a limitation. The intentional part: each of these figures
operates in a register where direct speech is a form of exposure,
and the novel tried to honour that. Guyon's ideas were dangerous
when stated plainly. Mackandal's communications were dangerous
by definition. Payne learned early to qualify what she said before
it reached the wrong ears. Julius posts under a name that is not
his name.

The limitation: direct dialogue is difficult to do well, and
this novel does not always do it well. The scenes that most
need voices---Guyon and Fénelon in the room together, Mackandal
and the network he is building, Payne in the moment of being
told to soften her conclusion---are scenes where the prose
retreats into interiority when it should risk the exposure of
speech. This is a failure of nerve as much as method, and a
future draft should correct it.

What the novel does not attempt, and should eventually attempt,
is the scene in which two of these registers meet directly---
in which the computational signatures that distinguish each
character become audible as friction rather than simply as
contrast. That scene does not appear here because it was not
yet possible to write it. It may become possible.

The objects---the surface of inscription, the instrument of
transformation, the observational frame---are the novel's
real connective tissue, more so than any thematic argument.
They accumulate meaning across the four parts in ways that
the author did not fully anticipate when they were introduced
and does not fully understand now. This seems correct. The
objects are doing what objects do in the frameworks this novel
draws on: carrying more than was put into them, available to
readings that the person who made them did not intend.

The reader will find what the reader finds.

\bigskip

\noindent\textit{Flyxion}\\
\textit{Independent Researcher}

